\documentclass[12pt]{article}
\usepackage[utf8]{inputenc}
\usepackage{amsmath}
\usepackage{graphicx}
\usepackage{natbib}
\usepackage{hyperref}
\usepackage{lipsum} % For dummy text, you can remove this in your final document

\title{Impact of Microplastics on Soil Health in Agricultural Ecosystems}
\author{Your Name}
\date{\today}

\begin{document}

\maketitle

\begin{abstract}
Microplastics, defined as plastic particles smaller than 5 mm, have become a significant environmental concern. This paper reviews the mechanisms of microplastic contamination in agricultural soils, analyzes their ecological consequences, and discusses potential mitigation strategies.
\end{abstract}

\section{Introduction}
Microplastics have gained attention due to their persistence and widespread presence in the environment. Initially studied in marine ecosystems, their impact on terrestrial environments, particularly in agricultural soils, is increasingly recognized. This paper aims to explore the sources and pathways of microplastic contamination in soils, assess their ecological consequences on soil health and agricultural productivity, and propose mitigation strategies to address this environmental challenge.

\section{Mechanisms of Microplastic Contamination}
Microplastics enter agricultural soils through various pathways:

\begin{itemize}
    \item Direct application of plastic mulches
    \item Sewage sludge application as fertilizer
    \item Atmospheric deposition
    \item Irrigation with contaminated water
\end{itemize}

These pathways highlight the multiple sources through which microplastics accumulate in agricultural ecosystems.

\section{Ecological Consequences}
The presence of microplastics in agricultural soils can lead to several ecological consequences:

\begin{itemize}
    \item Alterations in soil physical properties
    \item Adsorption of agrochemicals and pollutants
    \item Biological impacts on soil organisms
\end{itemize}

These consequences can affect soil fertility, nutrient cycling, and overall ecosystem health.

\section{Mitigation Strategies}
Addressing microplastic contamination in agricultural soils requires integrated mitigation strategies:

\begin{itemize}
    \item Regulation and policy measures
    \item Promotion of biodegradable mulches
    \item Improved plastic waste management
    \item Development of soil remediation techniques
\end{itemize}

Implementing these strategies can help minimize microplastic inputs and mitigate their adverse effects on soil health.

\section{Conclusion}
Microplastic contamination poses a significant threat to soil health in agricultural ecosystems, with potential implications for food security and environmental sustainability. Effective management strategies, supported by scientific research and policy interventions, are essential to mitigate these impacts and ensure the long-term health of agricultural soils.

\section*{References}
\bibliographystyle{plainnat}
\bibliography{prompt2_35}

\end{document}
